% !TeX root = ../../book.tex
\section{有限集}
\secbegin{secFiniteSets}

正如标题所示，本节主要探索有关有限集的属性，为此我们必须首先定义`有限'的含义。当我们看到一个有限集时，我们当然知道它是有限的 --- 例如：
\begin{itemize}
\item 集合 $\{ \text{红色}, \text{橙色}, \text{黄色}, \text{绿色}, \text{蓝色}, \text{紫色} \}$ 是有限的。
\item 集合 $[0,1]$ 无限的，但它有有限的长度。
\item 集合 $[0,\infty)$ 是无限的并具有无限的长度。
\item 集合 $\mathcal{P}(\mathbb{N})$ 是无限的，但没有所谓`长度'的概念。
\item 空集 $\varnothing$ 是有限的。
\end{itemize}

如果我们要给`有限集'下个定义，我们必须首先弄清楚上面的有限集有什么共同点而无限集却没有。

很难不精确地定义`有限'。首次尝试定义可能类似于以下内容：
\begin{center}
\textit{如果集合 $X$ 中的元素不能永远持续下去，则说集合 $X$ 是有限的。}
\end{center}
这是很好的直觉，但作为数学定义还不够好，因为`持续'和`永远'不是精确的术语（除非它们本身被定义）。因此，让我们尝试使其更加精确：
\begin{center}
\textit{如果集合 $X$ 中的元素能够逐一列出，由头有尾，则说集合 $X$ 是有限的。}
\end{center}
这个定义要好一些，但仍然不完全精确 --- `逐一列出'的定义并不完全清除。但我们可以精确地说：逐一列出 $X$ 的元素意味着我们指定`第一个元素'、`第二个元素'、`第三个元素'，依此类推。说这个列表有一个终点意味着我们最终到达`第 $n$ 个元素'，并且对于某个 $n \in \mathbb{N}$，不存在 `第 $(n +1) $ 个元素'。换句话说，对于某个自然数 $n$，我们将 $X$ 的元素与从 $1$ 到 $n$ 的自然数配对。

回想一下，对于每个 $n \in \mathbb{N}$，从 $1$ 到 $n$ 的自然数集都有自己的符号：

\RSdefBracketN*

由于`配对'实际上意味着`找到双射'，我们现在可以定义有限集的含义了。

\begin{definition}
\label{defFiniteSet}
\index{finite set}
\index{infinite set}
\index{enumeration!of a finite set}
如果对于某个 $n \in \mathbb{N}$ 存在双射 $f : [n] \to X$，则说集合 $X$ 是\textbf{有限的}。函数 $f$ 被称为 $X$ 的\textbf{枚举}。如果 $X$ 不是有限的，则说它是\textbf{无限的}。
\end{definition}

这个定义给出了以下证明集合有限的策略。

\begin{strategy}[证明一个集合是有限集]
为了证明集合 $X$ 是有限的，只需对于某个 $n \in \mathbb{N}$ 找到双射 $[n] \to X$ 即可。
\end{strategy}

\begin{example}
\label{exSomeFiniteSets}
设 $X = \{ \text{红色}, \text{橙色}, \text{黄色}, \text{绿色}, \text{蓝色}, \text{紫色} \}$。上文中我们指出 $X$ 是有限的；现在我们来证明它。定义 $f : [6] \to X$ 使得
\[ \begin{matrix}
f(1) = \text{红色} & f(2) = \text{橙色} & f(3) = \text{黄色} \\
f(4) = \text{绿色} & f(5) = \text{蓝色} & f(6) = \text{紫色}
\end{matrix} \]
函数 $f$ 显然是一个双射，因为对于某个 $k \in [6]$, $X$ 的每个元素都可以唯一地表示为 $f(k)$。所以$X$是有限的。
\end{example}

\begin{exercise}
\label{exBracketNIsFinite}
证明对于每个 $n \in \mathbb{N}$, $[n]$ 是有限的。
\end{exercise}

请注意，\Cref{exBracketNIsFinite} 特别指出 $\varnothing$ 是有限的，因为 $\varnothing = [0]$。

\subsection*{有限集的大小}

有时仅仅知道集合是有限的会很有用，但知道它有多少个元素会更有用。这个数量称为集合的\textit{大小}。直观上，集合的大小应该是其元素列表的长度，但为了明确这一点，我们首先需要知道列表中元素的数量与我们选择列出它们的方式无关。

有限集 $X$ 的`元素列表'是 \Cref{defFiniteSet} 给出的双射 $[n] \to X$，并且 $n$ 是列表的长度，这意味着我们需要证明如果 $[m] \to X$ 和 $[n] \to X$ 是双射，则 $m=n$。这就是 \Cref{thmUniquenessOfSize}。

为了证明这一点，我们必须首先证明一些在证明中会用到的涉及函数的结论。

\begin{lemma}
\label{lemRemoveElementFromDomainAndCodomain}
设 $X$ 和 $Y$ 为集合，设 $f : X \to Y$ 为单射并设 $a \in X$。则对于所有 $x \in X \setminus \{ a \}$ 存在由 $f^{\vee}(x) = f(x)$ 定义的单射 $f^{\vee} : X \setminus \{ a \} \to Y \setminus \{ f(a) \}$。此外，如果 $f$ 为双射 $X \to Y$，则 $f^{\vee}$ 为双射 $X \setminus \{ a \} \to Y \setminus \{ f(a) \}$。
\end{lemma}

\begin{cproof}
注意由于 $f$ 的单射性，$f^{\vee}$ 是明确定义的：事实上，如果对某个 $x \in X \setminus \{ a \}$ 我们有 $f(x) = f(a)$，那么这意味着 $x=a \in X \setminus \{ a \} $，这是一个矛盾。

要证明 $f^{\vee}$ 为单射，设 $x,y \in X \setminus \{ a \}$ 且 $f^{\vee}(x) = f^{\vee}(y)$。则 $f(x) = f(y)$, 所以根据 $f$ 的单射性可得 $x=y$。

现在假设 $f$ 为双射。则 $f$ 具有逆 $g : Y \to X$，其本身也是一个双射（因此也是一个单射）。对于所有 $y \in Y$，对 $g$ 重复上面的论证得到一个由 $g^{\vee}(y) = g(y)$ 定义的单射 $g^{\vee} : Y \setminus \{ f(a) \} \to X \setminus \{ a \}$。那么对于所有 $x \in X \setminus \{ a \}$ 和所有 $y \in Y \setminus \{ f(a) \}$ 我们有
\[ g^{\vee}(f^{\vee}(x)) = g(f(x)) = x \quad \text{and} \quad f^{\vee}(g^{\vee}(y)) = f(g(y)) \]
所以 $g^{\vee}$ 是 $f^{\vee}$ 的逆。因此如果 $f$ 为双射，则 $f^{\vee}$ 也为双射。
\end{cproof}

\Cref{lemRemoveElementFromDomainAndCodomain} 带来了下一个有用的结论，我们将在证明每个集合都有唯一的`大小'时使用它。

\begin{lemma}
\label{lemRemoveElementFromSet}
设 $X$ 为非空集合。对于所有 $a, b \in X$ 存在双射 $X \setminus \{ a \} \to X \setminus \{ b \}$。
\end{lemma}

\begin{cproof}
定义 $f : X \to X$ 为函数，交换 $a$ 和 $b$ 并保持 $X$ 的所有其他元素不变 --- 也就是说，我们定义
\[ f(x) = \begin{cases} b & \text{如果 } x = a \\ a & \text{如果 } x = b \\ x & \text{如果 } x \not\in \{ a, b \} \end{cases} \]
注意，对于所有 $x \in X$, $f(f(x)) = x$，所以 $f$ 是一个双射 --- 它有自己的逆！因为 $f(a) = b$，根据 \Cref{lemRemoveElementFromDomainAndCodomain} 可得对于所有 $x \in X \setminus \{ a \}$ 由 $f^{\vee}(x) = f(x)$ 定义的函数 $f^{\vee} : X \setminus \{ a \} \to X \setminus \{ b \}$ 是一个双射。
\end{cproof}

\begin{theorem}
\label{thmJectionsAndSizeOfNaturalNumbers}
设 $m,n \in \mathbb{N}$。
\begin{enumerate}[(a)]
\item 如果存在单射 $f : [m] \to [n]$，则 $m \le n$。
\item 如果存在满射 $g : [m] \to [n]$，则 $m \ge n$。
\item 如果存在双射 $h : [m] \to [n]$，则 $m=n$。
\end{enumerate}
\end{theorem}

\begin{cproof}[of (a)]
对于给定 $m \in \mathbb{N}$，设 $p(m)$ 为断言，对于所有 $n \in \mathbb{N}$，如果存在单射 $[m] \to [n]$，则 $m \le n$。我们用归纳法证明对于所有 $m \in \mathbb{N}$,  $p(m)$ 都为真。

\begin{itemize}
\item (\textbf{基本情况}) 我们需要证明，对于所有 $n \in \mathbb{N}$ 如果存在单射 $[0] \to [n]$，则 $0 \le n$。 这显然成立，因为对于所有 $n \in \mathbb{N}$, $0 \le n$。
\item (\textbf{归纳步骤}) 给定 $m \in \mathbb{N}$ 并假设对于所有 $n \in \mathbb{N}$，如果存在单射 $[m] \to [n]$，则 $m \le n$。

设 $n \in \mathbb{N}$ 并假设存在单射 $f : [m+1] \to [n]$。我们需要证明 $m+1 \le n$。

首先注意 $n \ge 1$。事实上，因为 $m+1 \ge 1$，我们有 $1 \in [m+1]$，所以 $f(1) \in [n]$。这意味着 $[n]$ 非空，所以 $n \ge 1$。特别地，$n-1 \in \mathbb{N}$ 所以集合 $[n-1]$ 是明确定义的。只需证明 $m \le n-1$ 即可。

设 $a = f(m+1) \in [n]$。则：
\begin{itemize}
\item 根据 \Cref{lemRemoveElementFromDomainAndCodomain}，存在单射 $f^{\vee} : [m] = [m+1] \setminus \{ m+1 \} \to [n] \setminus \{ a \}$。
\item 根据 \Cref{lemRemoveElementFromSet} 存在双射 $g : [n] \setminus \{ a \} \to [n] \setminus \{ n \} = [n-1]$。
\end{itemize}

组合这两个函数可得单射 $g \circ f^{\vee} : [m] \to [n-1]$。根据归纳假设有 $m \le n-1$，所以 $m+1 \le n$。
\end{itemize}
由归纳法可得出结论。
\end{cproof}

\begin{exercise}
证明 \Cref{thmJectionsAndSizeOfNaturalNumbers} 的 (b) 和 (c) 结论。
\hintlabel{exFinishProofOfJectionsAndSizeOfNaturalNumbers}{%
(b) 可以通过与 (a) 非常类似的的归纳法证明。在归纳步骤中，给定一个满射 $g : [m+1] \to [n]$，必有 $n \ge 1$，对于某个合适的 $a \in [m+1]$，构造一个满射 $g^- : [m+1] \setminus \{ a \} \to [n-1]$。然后调用 \Cref{lemRemoveElementFromSet}，使用 $[m] = [m+1] \setminus \{m+1\}$。
}
\end{exercise}
唷！蛮好玩的。证明了这些技术结论后，我们现在可以证明明确定义的有限集的大小所需的定理。

\begin{theorem}[大小唯一性]
\label{thmUniquenessOfSize}
设 $X$ 为有限集并设 $f : [m] \to X$ 和 $g : [n] \to X$ 为 $X$ 的枚举，其中 $m,n \in \mathbb{N}$。则 $m=n$。
\end{theorem}

\begin{cproof}
因为 $f : [m] \to X$ 和 $g : [n] \to X$ 为双射，根据 \Cref{exCompositeOfBijectionsIsBijection,exInverseBijection}，函数 $g^{-1} \circ f : [m] \to [n]$ 也为双射。因此根据 \Cref{thmJectionsAndSizeOfNaturalNumbers}(c)，$m=n$。
\end{cproof}

正如我们上面提到的，\Cref{thmUniquenessOfSize} 告诉我们，列出（枚举）有限集元素的任意两种方法都会产生相同数量的元素。我们现在可以做出以下定义。

\begin{definition}
\label{defSize}
\index{size}
\index{cardinality!of a finite set}
设 $X$ 为有限集。$X$ 的\textbf{大小}（或 \textbf{基数}），写做 $|X|$，是一个唯一的自然数 $n$，满足存在双射 $[n] \to X$。
\end{definition}

\begin{example}
\Cref{exSomeFiniteSets} 给出 $|\{ \text{红色}, \text{橙色}, \text{黄色}, \text{绿色}, \text{蓝色}, \text{紫色} \}| = 6$，并给出证明证明是正确的，\Cref{exBracketNIsFinite} 给出对于所有 $n \in \mathbb{N}$, $|[n]| = n$；特别的，$|\varnothing| = 0$。
\end{example}

\begin{example}
\label{exBijectionFromIntegersFromMinusNToNToBracketTwoNPlusOne}
给定 $n \in \mathbb{N}$ 并设 $X = \{ a \in \mathbb{Z} \mid -n \le a \le n \}$。存在定义为 $f(k) = k-n-1$ 的双射 $f : [2n+1] \to X$。事实上：
\begin{itemize}
\item \textbf{$f$ 是明确定义的。}
%% BEGIN EXTRACT (xtrIntroducingVariableExampleTwo) %%
我们需要证明对于所有 $k \in [2n+1]$, $f(k) \in X$。给定 $k \in [2n+1]$，我们有 $1 \le k \le 2n+1$，所以
\[ -n = 1-(n+1) \le ~ \underbrace{k-(n+1)}_{=f(k)} ~  \le (2n+1) - (n+1) = n \]
因此 $f(k) \in X$。
%% END EXTRACT %%
\item \textbf{$f$ 是单射} 设 $k, \ell \in [2n+1]$ 并假设 $f(k) = f(\ell)$。那么 $k - n - 1 = \ell - n - 1$，所以 $k = \ell$。
\item \textbf{$f$ 是满射} 设 $a \in X$ 并定义 $k = a+n+1$。那么
\[ 1 = (-n) + n + 1 \le ~ \underbrace{a + n + 1}_{= k} ~ \le n + n + 1 = 2n+1 \]
所以 $k \in [2n+1]$，此外 $f(k) = (a+n+1)-n-1 = a$。
\end{itemize}
所以 $f$ 为双射，根据 \Cref{defSize} 我们有 $|X| = 2n+1$。
\end{example}

\begin{exercise}
\label{exAddRemoveElementsOfFiniteSets}
设 $X$ 为有限集，其中 $|X| = n > 1$。设 $a \in X$ 并设 $b \not \in X$。证明
\begin{enumerate}[(a)]
\item $X \setminus \{ a \}$ 为有限集，且 $|X \setminus \{ a \}| = n-1$；
\item $X \cup \{ b \}$ 为有限集，且 $|X \cup \{ b \}| = n+1$。
\end{enumerate}
确定在证明中何处使用 $a \in X$ 和 $b \not \in X$ 假设。
\end{exercise}

\subsection*{比较有限集的大小}

我们在 \Cref{secInjectionsSurjections} 开头部分使用点和星来引出单射和满射函数的定义时，提出了以下直觉：
\begin{itemize} 
\item 如果存在单射 $f : X \to Y$，则 $X$ `最多与 $Y$ 的元素一样多'；
\item 如果存在满射 $g : X \to Y$，则 $X$ `至少与 $Y$ 的元素一样多'。
\end{itemize}

现在是时候来证明至少当 $X$ 和 $Y$ 为有限集时上面的结论成立。下面的定理是对 \Cref{thmJectionsAndSizeOfNaturalNumbers} 的泛化。

\begin{theorem}
\label{thmFiniteSetsAndJections}
设 $X$ 和 $Y$ 为集合。
\begin{enumerate}[(a)]
\item 如果 $Y$ 为有限集且存在单射 $f : X \to Y$，则 $X$ 为有限集且 $|X| \le |Y|$。
\item 如果 $X$ 为有限集且存在满射 $f : X \to Y$，则 $Y$ 为有限集且 $|X| \ge |Y|$。
\item 如果 $X$ 或 $Y$ 其中之一为有限集且存在双射 $f : X \to Y$，则 $X$ 和 $Y$ 都是有限集且 $|X| = |Y|$。
\end{enumerate}
\end{theorem}

\begin{cproof}[of (a)]
\item 我们通过归纳法来证明，对于所有 $n \in \mathbb{N}$，如果 $Y$ 是大小为 $n$ 的有限集且存在单射 $f : X \to Y$，则 $X$ 为有限集且 $|X| \le n$。

\begin{itemize}
\item (\textbf{基本情况}) 设 $Y$ 是大小为 $0$ 的有限集 --- 也就是说，$Y$ 为空集。假设存在单射 $f : X \to Y$。如果 $X$ 非空，则存在元素 $a \in X$，使得 $f(a) \in Y$。这与 $Y$ 为空集相矛盾，所以 $X$ 一定为空。因此 $|X| = 0 \le 0$。

\item (\textbf{归纳步骤}) 给定 $n \in \mathbb{N}$ 并假设如果 $Z$ 是大小为 $n$ 的有限集，且存在单射 $f : X \to Z$，则 $X$ 为有限集且 $|X| \le n$。

给定大小为 $n+1$ 的有限集 $Y$ 和单射 $f : X \to Y$。我们需要证明 $X$ 为有限集且 $|X| \le n+1$。

如果 $X$ 为空集，则 $|X| = 0 \le n+1$。假设 $X$ 非空，给定元素 $a \in X$。

根据 \Cref{lemRemoveElementFromDomainAndCodomain}，存在单射 $f^{\vee} : X \setminus \{ a \} \to Y \setminus \{ f(a) \}$。此外，根据 \Cref{exAddRemoveElementsOfFiniteSets}，$Y \setminus \{ f(a) \}$ 为有限集且 $|Y \setminus \{ f(a) \}| = (n+1) - 1 = n$。

应用归纳假设 $Z = Y \setminus \{ f(a) \}$ 可知 $X \setminus \{ a \}$ 为有限集且 $|X \setminus \{ a \}| \le n$。而根据 \Cref{exAddRemoveElementsOfFiniteSets}，$|X \setminus \{ a \}| = |X| - 1$，所以 $|X| \le n+1$。
\end{itemize}

根据归纳法得出结论。
\end{cproof}

\begin{exercise}
证明 \Cref{thmFiniteSetsAndJections} 中的 (b) 和 (c)。
\end{exercise}

\Cref{thmFiniteSetsAndJections} 给出了如下比较两个有限集大小的策略：
\begin{strategy}[比较有限集大小]
\label{strComparingSizesOfFiniteSets}
设 $X$ 和 $Y$ 为有限集。
\begin{enumerate}[(a)]
\item 要想证明 $|X| \le |Y|$，只需找到 $X \to Y$ 的单射即可。
\item 要想证明 $|X| \ge |Y|$，只需找到 $X \to Y$ 的满射即可。
\item \label{strBijectiveProof} 要想证明 $|X| = |Y|$，只需找到 $X \to Y$ 的双射即可。
\end{enumerate}
策略 \ref{strBijectiveProof} 通常被称为\textbf{双射证明}。
\end{strategy}

\subsection*{有限集的封闭性}

我们现在使用 \Cref{strComparingSizesOfFiniteSets} 来证明有限集的\textit{封闭性} --- 也就是说，我们可以对有限集执行操作以确保结果依然有限。

\begin{exercise}
设 $X$ 为有限集。证明每个子集 $U \subseteq X$ 为有限集且 $|U| \le |X|$。
\hintlabel{exSubsetOfFiniteSetIsFinite}{%
给定 $U \subseteq X$，找到 $U \to X$ 的单射并应用 \Cref{thmFiniteSetsAndJections}(a)。
}
\end{exercise}

\begin{exercise}
设 $X$ 和 $Y$ 为有限集。证明 $X \cap Y$ 为有限集。
\hintlabel{exIntersectionOfFiniteSetsIsFinite}{%
回忆一下对于所有集合 $X$ 和 $Y$, $X \cap Y \subseteq X$。
}
\end{exercise}

\begin{proposition}
\label{propUnionOfFiniteSetsIsFinite}
设 $X$ 和 $Y$ 为有限集，则 $X \cup Y$ 为有限集，并且
\[ |X \cup Y| = |X| + |Y| - |X \cap Y| \]
\end{proposition}

\begin{cproof}
我们先在 $X$ 和 $Y$ 不相交的情况下证明这一点。更一般的情况，即不假设 $X$ 和 $Y$ 不相交将在 \Cref{exSizeOfUnion} 中证明。

设 $m = |X|$ 且 $n=|Y|$，并设 $f : [m] \to X$ 且 $g : [n] \to Y$ 为双射。
 
因为 $X$ 和 $Y$ 不相交，可得 $X \cap Y = \varnothing$。定义 $h : [m+n] \to X \cup Y$ 如下；给定 $k \in [m+n]$，设
\[ h(k) = \begin{cases} f(k) & \text{if } k \le m \\ g(k-m) & \text{if } k > m \end{cases} \]
注意，$h$ 是明确定义的：情况 $k \le m$ 和 $k > m$ 是互斥的，他们覆盖了所有可能情况，且对于所有 $m+1 \le k \le n$, $k-m \in [n]$，所以 $g(k-m)$ 有定义。然后可以直接检查 $h$ 具有逆 $h^{-1} : X \cup Y \to [m+n]$ 定义为
\[ h^{-1}(z) = \begin{cases} f^{-1}(z) & \text{if } z \in X \\ g^{-1}(z)+m & \text{if } z \in Y \end{cases} \]
$h^{-1}$ 的明确定义从根本上依赖于假设 $X \cap Y = \varnothing$，因为这可以确保 $x \in X$ 和 $x \in Y$ 是互斥的。

%% BEGIN EXTRACT (xtrConclusionExampleTwo) %%
因此 $|X \cup Y| = m+n = |X| + |Y|$，因为 $|X \cap Y| = 0$。
%% END EXTRACT %%
\end{cproof}

\begin{exercise}
\label{exSizeOfUnion}
按照如下步骤完成 \Cref{propUnionOfFiniteSetsIsFinite} 的证明：
\begin{enumerate}[(a)]
\item 给定集合 $A$ 和 $B$，证明集合 $A \times \{ 0 \}$ 和 $B \times \{ 1 \}$ 不相交，并找到 $A \to A \times \{ 0 \}$ 和 $B \to B \times \{ 1 \}$ 的双射。写做 $A \sqcup B$\nindex{unionDisjoint}{$\sqcup$}{disjoint union} \inlatex{sqcup}\lindexmmc{sqcup}{$\sqcup$} 表示集合 $(A \times \{ 0 \}) \cup (B \times \{ 1 \})$。集合 $A \sqcup B$ 称为 $A$ 和 $B$ 的\textbf{不相交并集}\index{disjoint union}。
\item 证明，如果 $A$ 和 $B$ 为有限集，则 $A \sqcup B$ 为有限集，且
\[ |A \sqcup B| = |A| + |B| \]
\item 设 $X$ 和 $Y$ 为集合。找到双射
\[ (X \cup Y) \sqcup (X \cap Y) \to X \sqcup Y \]
\item 完成 \Cref{propUnionOfFiniteSetsIsFinite} 的证明 --- 即，证明如果 $X$ 和 $Y$ 为有限集，不需要相交与否，$X \cup Y$ 为有限集，且
\[ |X \cup Y| = |X| + |Y| - |X \cap Y| \]
\end{enumerate}
\end{exercise}

\begin{exercise}
设 $X$ 为有限集，并设 $U \subseteq X$。证明 $X \setminus U$ 为有限集，且 $|X \setminus U| = |X| - |U|$。
\hintlabel{exSizeOfRelativeComplement}{%
将 \Cref{propUnionOfFiniteSetsIsFinite} 与 $Y=U$ 一起应用。
}
\end{exercise}

\begin{exercise}
\label{exSizeOfCartesianProduct}
设 $m,n \in \mathbb{N}$。证明 $|[m] \times [n]| = mn$。
\hintlabel{exBijectionBetweenProductOfBracketSets}{%
在 $n \in \mathbb{N}$ 上用归纳法证明，对于所有 $m,n \in \mathbb{N}$，存在 $[mn] \to [m] \times [n]$ 的双射。
}
\end{exercise}

\begin{proposition}
\label{propProductOfFiniteSetsIsFinite}
设 $X$ 和 $Y$ 为有限集。那么 $X \times Y$ 为有限集，且
\[ |X \times Y| = |X| \cdot |Y| \]
\end{proposition}
\begin{cproof}
设 $X$ 和 $Y$ 为有限集，设 $m=|X|$ 和 $n=|Y|$，并设 $f : [m] \to X$ 和 $g : [n] \to Y$ 为双射。定义函数 $h : [m] \times [n] \to X \times Y$ 为
\[ h(k,\ell) = (f(k),g(\ell)) \]
对于每个 $k \in [m]$ 和 $\ell \in [n]$。很容易发现这是一个双射，其逆为
\[ h^{-1}(x,y) = (f^{-1}(x), g^{-1}(y)) \]
对于所有 $x \in X$ 和 $y \in Y$。根据 \Cref{exBijectionBetweenProductOfBracketSets} 存在双射 $p : [mn] \to [m] \times [n]$，并根据 \Cref{exCompositeOfBijectionsIsBijection}，组合 $h \circ p : [mn] \to X \times Y$ 也为双射。因此 $|X \times Y| = mn$。
\end{cproof}

总之，我们证明了通过取子集、成对并集、成对交集、成对笛卡尔积和相对补集可以保持有限性。
\subsection*{无限集}

最后，我们通过证明并非所有集合都是有限的来结束本节 --- 具体来说，我们将证明 $\mathbb{N}$ 是无限的。\textit{直觉上}这似乎显而易见：$\mathbb{N}$ 理所\textit{当然}是无限的！但在数学实践中，这是不够的：我们需要使用`无限'的定义来证明 $\mathbb{N}$ 是无限的。也就是说，我们需要证明对于任意 $n \in \mathbb{N}$ 不存在双射 $[n] \to \mathbb{N}$。我们将在下面的证明中使用 \Cref{lemFinSetGreatestElement}。

\begin{lemma}
\label{lemFinSetGreatestElement}
自然数的每个非空有限集都有一个最大元素。
\end{lemma}
\begin{cproof}
我们用归纳法证明在 $n \ge 1$ 时，每个大小为 $n$ 的子集 $U \subseteq \mathbb{N}$ 都有最大元素。
\begin{itemize}
\item (\textbf{基本情况}) 让 $U \subseteq \mathbb{N}$ 且 $|U|=1$。则对于某个  $m \in \mathbb{N}$, $U = \{ m \}$。因为 $m$ 是 $U$ 的唯一元素，它当然是 $U$ 的最大元素！

\item (\textbf{归纳步骤}) 给定 $n \ge 1$ 并假设每个大小为 $n$ 的自然数集都有最大元素（\textbf{IH}）。

设 $U \subseteq \mathbb{N}$ 且 $|U| = n+1$。我们希望证明 $U$ 存在最大元素。

因为 $|U| = n+1$，对于不同自然数 $m_k$ 我们可以写做 $U = \{ m_1, m_2, \dots, m_n, m_{n+1} \}$。那么根据 \Cref{exAddRemoveElementsOfFiniteSets}，$|U \setminus \{ m_{n+1} \}| = n$，根据归纳假设，$U \setminus \{ m_{n+1} \}$ 存在最大元素，假设为 $m_k$。则：
\begin{itemize}
\item 如果 $m_k < m_{n+1}$，则 $m_{n+1}$ 为 $U$ 的最大元素。
\item 如果 $m_k > m_{n+1}$，则 $m_k$ 为 $U$ 的最大元素。
\end{itemize}
无论哪种情况，$U$ 都有最大元素。归纳步骤完毕。
\end{itemize}
\end{cproof}

\begin{theorem}
\label{thmNIsInfinite}
集合 $\mathbb{N}$ 是无限的。
\end{theorem}

\begin{cproof}
%% BEGIN EXTRACT (xtrContradictionExample) %%
我们通过反证法来证明。假设 $\mathbb{N}$ 是有限的。那么对于某个 $n \in \mathbb{N}$, $|\mathbb{N}| = n$，因为 $\mathbb{N}$ 要么是空集（显然不可能，因为 $0 \in \mathbb{N}$），要么根据 \Cref{lemFinSetGreatestElement} 存在最大元素 $g$。然而由于每个自然数都有后继，所以 $g+1 \in \mathbb{N}$，而 $g+1 > g$，这与 $g$ 为最大值矛盾。因此 $\mathbb{N}$ 是无限的。
%% END EXTRACT %%
\end{cproof}